\section*{Seznámení s tématem}
\addcontentsline{toc}{section}{Seznámení s tématem}
Zdraví člověka je jedním z nejdůležitějších témat dnešní doby. Tuto skutečnost lze pozorovat například na rozpočtu vlády ČR na rok 2025, kde Ministerstvu zdravotnictví byl zvýšen rozpočet, o proti roku 2024 o zhruba 53,9 \% \parencite{prcRozpocet2025}. Jeden z důvodů předčasného úmrtí člověka je, mimo kouření a nezdravé stravování, znečištění ovzduší. Znečištění ovzduší stojí za miliony předčasných úmrtí ročně a je podle WHO největší enviromentální hrozbou pro zdraví člověka \parencite{worldhealthorganizationWHOGlobalAir2021}. 

Znečistěné ovzduší se netýká pouze vnějších prostor, ale také těch vnitřních. Ovzduší ve vniřních prostorech by se dalo považovat, v jistých ohledech, za důležitější než to vnější. To z toho důvodu, že většina aktivity moderního člověka se odehrává právě v nějakém uzavřeném prostoru. Ať už se jedná o kancelář, školu, restauraci, bar, obchod nebo dokonce dům, je zřejmé, že aktivita člověka probíhá právě v uzavřených prostorech. V těchto prostorech může docházet k environmentálním jevům, které výrazně ovlivňují fyziologii člověka a jeho produktivitu a komfort \parencite{dengImpactIndoorAir2024}.